% ============================================================
\section{Problema 3: Configuración de topología del VLSM II}
% ============================================================

\begin{tcolorbox}[enunciado]
En la asignación se le adjunta un \texttt{.pkt} en donde solamente tendrá que hacer las configuraciones correspondientes.

Usando el VLSM que en el anterior ejercicio calculó, entonces configure la topología considerando estos puntos:
\begin{itemize}
    \item Todos los dispositivos intermedios deben tener hostname, el cual es el que viene en la etiqueta del dispositivo por defecto, por ejemplo cuando se arrastra un router dice router1.
    \item Para todos los Switches use contraseña tipo 7 y Telnet.
    \item Para todos los Router deben usar contraseña tipo 5 y configurado SSH. 
    \item Para el dominio de ssh use \texttt{ciencias.unam.mx}.
    \item Todos los dispositivos intermedios deben tener contraseña de enable, en este caso para switches debe ser 7 y routers tipo 5. 
    \item Debe ir enrutado con RIP.
    \item Los routers (EBIO, RFIS, RORI, RMAT, RAMO, RPO, CRouter/taller, RYELIZ, RTLATIPAC, RTLAH, C3Router/EXBIB) deben dar DHCP \textbf{solamente a los dispositivos finales excepto los servidores.} 
    
    \textbf{Observación:} Ponga atención en qué segmento esté dando DHCP ya que no es a todos, sino como corresponde al departamento.
    \item Debe configurarse bien el servicio de SMTP y WEB. En el caso de servicio de SMTP todas las computadoras y laptops deben tener un usuario y contraseña, además de ello, debe ser dominio \texttt{ciencias.unam.mx}. 
    \item En el FTP configure una cuenta por cada departamento e intente subir un archivo en texto en claro.
    \item Para servicio HTTP sería \texttt{www.redfciencias.com} y para HTTPS \texttt{web.fciencias.unam.mx}.
    \item Como prueba que sí funciona la configuración, es hacer ping.
    \item Personalice las dos páginas usando HTML y CSS. Investigue como hacerlo.
\end{itemize}

Conteste las siguientes preguntas:
\begin{itemize}
    \item ¿Cuántos saltos muestra la tabla de enrutamiento si estamos desde el edificio de biología hacia el segmento de edificio del Tlahuizcalpan?
    \item La laptop0 quiere hacer un ping a la laptop6 ¿Es la misma IP que se quiere comunicar de la laptop0 y laptop6? Hint: Revise las IP asignadas por el DHCP y también recuerde el protocolo ARP.
\end{itemize}
\end{tcolorbox}

\subsection*{Solución}
% Escribe aquí tu solución para el Problema 3...